\documentclass[../DoAn.tex]{subfiles}
\begin{document}
\begin{center}
    \Large{\textbf{DANH MỤC THUẬT NGỮ}}\\
\end{center}
\vspace{0.5cm}

\begin{longtable}{|p{3.8cm}|p{10cm}|}
\hline
\multicolumn{1}{|c|}{\textbf{Thuật ngữ}} & \multicolumn{1}{c|}{\textbf{Ý nghĩa}} \\
\hline
\endfirsthead
\hline
\multicolumn{1}{|c|}{\textbf{Thuật ngữ}} & \multicolumn{1}{c|}{\textbf{Ý nghĩa}} \\
\hline
\endhead
\endfoot
\endlastfoot
Chấm công hộ & Một người chấm công thay cho người khác \\
\hline
Lượng tử hóa & Biểu diễn trọng số và phép tính của mô hình bằng số nguyên ít bit để giảm bộ nhớ và thời gian suy luận \\
\hline
Lượt chấm & Một lần ghi nhận một người tại máy chấm công \\
\hline
Mẫu khuôn mặt & Vector đặc trưng của một người được lưu để so khớp \\
\hline
Ngày công & Kết quả quy đổi các lượt chấm trong một ngày theo ca làm \\
\hline
Tấn công trình diện & Đưa ảnh, video hoặc vật giả trước cảm biến để đánh lừa hệ thống sinh trắc học \\
\hline
TinyML & Hướng chạy mô hình học máy trên vi điều khiển có tài nguyên hạn chế \\
\hline
Vector đặc trưng & Dãy số mô tả một khuôn mặt, do mô hình nhận diện tạo ra \\
\hline
\end{longtable}

\end{document}
